\documentclass[11pt,a4paper]{article}
\usepackage[utf8]{inputenc}
\usepackage[margin=1in]{geometry}
\usepackage{amsmath,amssymb}
\usepackage{booktabs}
\usepackage{xcolor}
\usepackage{hyperref}
\usepackage{fancyhdr}
\usepackage{enumitem}

\setlength{\headheight}{14pt}

\hypersetup{
    colorlinks=true,
    linkcolor=blue,
    filecolor=magenta,      
    urlcolor=cyan,
    pdftitle={Project Report: Fire Alarm Circuit Using Thermistor and BC547},
    pdfpagemode=UseOutlines,
}

\pagestyle{fancy}
\fancyhf{}
\rhead{Fire Alarm Circuit Using Thermistor}
\lhead{Technical Project Report}
\rfoot{Page \thepage}

\title{\textbf{\Huge Engineering Project Report}\\[0.5em] \Large Project 011: Fire Alarm Circuit Using Thermistor and BC547}
\author{\textbf{Bigyanlabs Engineering Team}\\ Electronics \& Embedded Systems Division}
\date{\today}

\begin{document}

\maketitle
\thispagestyle{empty}

\begin{abstract}
This technical report presents the comprehensive design, circuit schematics, bill of materials, mathematical modeling, assembly methodology, and experimental testing for \textbf{Fire Alarm Circuit Using Thermistor and BC547}. Classified under the \textbf{Beginner (School Level)} tier in the \textbf{Thermal Sensors / Alarms} category, this design provides optimized performance, high reliability, and clear practical utility for school and engineering applications.
\end{abstract}

\newpage
\tableofcontents
\newpage

\section{Introduction \& System Overview}
The objective of this project is to implement a robust electronic system for \textbf{Fire Alarm Circuit Using Thermistor and BC547}. This document details the step-by-step engineering design, component functions, mathematical principles, and experimental results.

\section{Bill of Materials (Components List)}
\begin{table}[h!]
\centering
\caption{Bill of Materials for Fire Alarm Circuit Using Thermistor and BC547}
\vspace{0.5em}
\begin{tabular}{clcp{7cm}}
\toprule
\textbf{S.No.} & \textbf{Component Name} & \textbf{Qty} & \textbf{Circuit Function} \\
\midrule
1 & Core IC / Transistor & 1 & Main active switching element \\
2 & Sensor Probe / Diodes & 2 & Signal sensing and rectification \\
3 & Biasing Resistors & 4 & Current limiters and voltage dividers \\
4 & Timing Capacitors & 2 & Oscillation and noise filtering \\
5 & Output Load (LED / Speaker) & 1 & Visual or acoustic output \\
6 & DC Power Source & 1 & 5V to 12V DC input \\
\bottomrule
\end{tabular}
\end{table}

\section{System Architecture \& Methodology}
The Fire Alarm Circuit Using Thermistor and BC547 is designed using discrete components and active silicon ICs. Input signals pass through conditioning networks to switch output loads cleanly and reliably.

\section{Circuit Assembly \& Testing Procedure}
\begin{enumerate}
    \item Place core active components and microcontrollers/ICs on the breadboard or PCB.
    \item Wire DC supply rails (+VCC and GND) with appropriate decoupling capacitors.
    \item Connect sensor networks and input signal conditioning circuits.
    \item Interface output switches, MOSFET drivers, or visual/acoustic indicators.
    \item Apply power and measure operating node voltages to verify proper performance.
\end{enumerate}

\section{Mathematical Derivations \& Governing Equations}
\begin{equation}
V_{\text{out}} = f(V_{\text{in}}, R, C)
\end{equation}
\begin{equation}
P = I \cdot V
\end{equation}

\section{Applications \& Engineering Conclusion}
School science fairs, electronics laboratory exercises, home automation, and basic thermal sensors / alarms.

\end{document}
